\section{Related work}
\label{sec:related}

\paragraph{Source ownership and attribution.}
Attribution evaluation across multiple candidate sources is an established measurement problem \citep{attributionbench2024}; source-aware factuality work targets cross-source conflation directly \citep{provenanceguard2026}; and scientific claim--evidence benchmarks test whether retrieved evidence supports or contradicts a claim \citep{claimbench2025}. Claim-source retrieval sharpens the same point from the retrieval side, since recovering the source a scientific statement actually came from is a different task from finding a source that would support it \citep{claimsource2026}. Together these establish claim correctness, source ownership, and semantic support as separate coordinates, which is the premise this paper builds on.

\paragraph{Provenance-sensitive authorization.}
A current controlled audit makes the authorization distinction especially direct: evidence can be relevant to an agent decision without being authorized to determine that decision, and target-specific tests can hold the task and proposition fixed while varying only source authority \citep{provsen2026}. The relevance-versus-authorized-evidence distinction is therefore already established. What remains open is the empirical question studied here: whether a protected conjunction of content, ownership, support, checker, influence, evaluator, access, and contamination prerequisites behaves differently from its strongest single component when both meet the same hostile battery.

\paragraph{Partial evidence and formal authorization.}
The donor frontier is stronger than provenance-sensitive factuality alone.  Partial Evidence Bench studies unsafe completeness and structured gap reporting when relevant evidence lies outside an authorized view \citep{partial2026}.  AuthGraph aligns execution provenance with an independently maintained authorization graph \citep{authgraph2026}; FAVA lowers natural-language constraints into evidence-backed permission graphs with deterministic authorization \citep{fava2026}; and CVA binds principal, request, context, and policy satisfaction in a cryptographically verifiable authorization object \citep{cva2026}.  P4 treats all four as donor-owned advances.  Its residual question is not whether evidence can be partial, provenance can be aligned with permission, or authorization can be formal.  It is whether the joint donor state is sufficient for the target \emph{scientific-promotion} terminal.  Section~\ref{sec:axis-theory} states the exact factorization condition: once a donor product exposes a target-sufficient state and the same terminal relation, it is extensionally equivalent.

\paragraph{Auditability, citation fidelity, and influence.}
Claim-level auditability work argues that research-agent outputs require persistent semantic provenance rather than only fluent text \citep{aar2026}. ProvenAI separates citation fidelity from whether cited evidence influenced the answer path \citep{provenai2026}. We adopt that separation: citation presence cannot substitute for evidence-use provenance.

\paragraph{Research-integrity provenance and declared behavior.}
INSPECT-AI and RIPE-KG represent the provenance of research-integrity assessments over clinical-trial publications, showing a current scientific setting in which the assessment process itself is an auditable object \citep{inspectai2026}. Behavioral Integrity Verification formalizes a complementary declared-versus-actual integrity problem for agent skills \citep{behavioralintegrity2026}, and certified speculative execution shows the same instinct applied to untrusted agent actions: let the work proceed, but hold its effects behind a check that must clear before they become authoritative \citep{certifiedspeculative2026}. These are parent-domain precedents for explicit provenance and contract checking.

\paragraph{Iterative verification and evidence escalation.}
FIRE iterates retrieval and verification rather than treating a single retrieval pass as final \citep{fire2024}. DeepSciVerify motivates escalation from an abstract-level judgment to fuller evidence when initial evidence is insufficient \citep{deepsciverify2026}. Stage-wise fact-checking benchmarks make the complementary measurement argument, that an end-to-end verdict hides which stage actually failed \citep{factarena2026}. The first two mechanisms inform frozen comparator proxies here, and the third informs the decision to report per-family and per-coordinate outcomes rather than one aggregate.

\paragraph{Evaluator integrity, benchmark auditing, and contamination.}
RewardHackingAgents treats evaluator manipulation and held-out leakage as an explicit benchmark dimension \citep{rewardhacking2026}. Search-time contamination work shows that browsing systems can retrieve public benchmark answers during inference \citep{stc2026}. BenchGuard and Automated Benchmark Auditing show that benchmark specifications, environments, ground truths, and grading logic can themselves be defective \citep{benchguard2026,autobenchmarkaudit2026}. The Holistic Agent Leaderboard further demonstrates the value of standardized harnessing and trace inspection, including direct observation of agents searching for benchmark material rather than solving the intended task \citep{hal2025}. These results motivate the protected evaluator identity, holdout custody, access telemetry, and independent audit used here.

The verification-axis analysis adds a distinction that benchmark auditing alone does not provide.  A task can be free of the registered nuisance shortcuts yet still fail to compare scientific judgement because the target terminal is absent from a system's interface; conversely, a perfectly expressive interface does not make a mixed information fibre identifiable.  Construction audit, information sufficiency, terminal attainability, and observed panel resolution must therefore be reported separately.

\paragraph{Abstention and scientific integrity.}
SciIntegrity-Bench constructs academic-integrity dilemmas in which honest acknowledgement of inability is the only correct outcome, while AgentAbstain evaluates paired should-act/should-abstain tasks in executable environments \citep{sciintegrity2026,agentabstain2026}. They establish abstention as an evaluation target distinct from generic task success. The object measured here is narrower: an undetermined or blocked terminal is authority-specific, produced by an unresolved or failed evidence, checker, or evaluator prerequisite rather than by a general judgement that the task should not be attempted.
